\chapter{Khi Thầy Cô Xin Lỗi}
\markboth{Khi Thầy Cô Xin Lỗi}{Khi Thầy Cô Xin Lỗi}

\begin{center}
	\textit{\color{bookprimary}``Xin lỗi đúng chỗ không làm người thầy nhỏ đi. Nó chỉ làm lớp học bớt giả hơn.''}
\end{center}

\ngancach

\begin{tcolorbox}[
	colback=bookprimary!6,
	colframe=bookprimary!35,
	arc=5pt,
	title={\bfseries Tại Sao Chương Này Quan Trọng?},
	fonttitle=\bfseries\color{bookprimary}
]
Nhiều học sinh lớn lên trong niềm tin rằng người lớn càng to tiếng càng đúng. Bởi vậy, một lời xin lỗi từ thầy cô thường gây chấn động hơn cả một bài giảng hay. Nhưng xin lỗi thế nào để không biến lớp học thành nơi ai cũng đòi được nuông chiều? Đó là điểm khó của chương này.
\end{tcolorbox}

\section{Xin Lỗi Nhưng Không Xuề Xòa}
\begin{truyen}{Xin Lỗi Nhưng Không Xuề Xòa}{A Firm Apology}
\chuhoa{S}{au khi kiểm tra lại bài, cô phát hiện mình đã chấm sót một ý và làm một bạn bị thấp hơn nửa điểm. Cô gọi bạn đó lên rồi nói trước lớp: ``Cô chấm thiếu. Cô xin lỗi em, và cô sửa điểm lại.''}

Bạn ấy hơi bất ngờ, nhưng cũng tranh thủ nói: ``Dạ vậy chắc mai mốt tụi em sai cô cũng nhẹ tay lại cho công bằng hen cô?''

\teacherpair{Cô mỉm cười: ``Cô xin lỗi vì cô sai. Còn các em làm sai thì vẫn phải chịu đúng hậu quả của phần sai đó. Công bằng không có nghĩa là ai sai cũng được xí xóa như nhau. Công bằng là ai sai phần nào thì sửa phần đó.''}{The teacher smiles: ``I apologize because I was wrong. When you are wrong, you still face the proper consequence of that mistake. Fairness does not mean everyone gets excused equally. Fairness means each wrong is corrected where it belongs.''}

Lớp im đi, kiểu im của những bạn vừa thấy một điều khó: xin lỗi thật ra không làm người ta yếu đi, nhưng cũng không làm kỷ luật mềm nhũn ra.

\textit{(The teacher admits grading unfairly, apologizes, and fixes it. But she refuses the student's attempt to turn that apology into a blanket reduction of standards.)}
\end{truyen}

\section{Xin Lỗi Vì Nặng Lời}
\begin{truyen}{Xin Lỗi Vì Nặng Lời}{Apology for Tone}
\chuhoa{C}{ó lần một thầy bực quá nên quát một học sinh trước lớp}. Hôm sau, thầy gọi em ở lại và nói ngắn: ``Hôm qua thầy nóng quá. Nội dung thầy nhắc em không sai, nhưng cách thầy nói là sai. Thầy xin lỗi đoạn đó.''

Bạn học sinh đứng khựng lại. Vì em không nghĩ người lớn lại tách bạch được giữa \textit{điều đúng} và \textit{cách sai} rõ đến vậy.

Thầy nói tiếp: ``Thầy vẫn giữ yêu cầu cũ: em phải sửa thái độ học tập. Nhưng thầy không có quyền dùng sự bực bội của mình để làm em bị nhục trước tập thể.''
\end{truyen}

\section{Xin Lỗi Trước Tập Thể}
\begin{truyen}{Xin Lỗi Trước Tập Thể}{Public Mistake, Public Repair}
\chuhoa{M}{ột buổi sinh hoạt}, cô chủ nhiệm trách oan cả lớp vì tưởng lớp giấu chuyện đánh nhau. Đến chiều, cô mới biết thông tin mình nghe chưa đầy đủ.

Sáng hôm sau cô đứng trước lớp: ``Hôm qua cô đã dùng một tin chưa kiểm chứng để quy kết cả lớp. Cô xin lỗi lớp.''

Một bạn nói nhỏ: ``Dạ tụi em bất ngờ lắm cô.''

\teacherpair{Bất ngờ là đúng. Vì nhiều người quen với việc người lớn chỉ xin lỗi riêng cho nhanh. Nhưng nếu lỗi của cô đã làm tổn thương tập thể trước đám đông, thì lời sửa cũng phải đứng trước tập thể.}{``Of course you are surprised. Many people are used to adults apologizing privately just to move on. But if my mistake hurt the group in public, the repair must also stand before the group.''}
\end{truyen}

\section{Xin Lỗi Không Phải Mặc Cả}
\begin{truyen}{Xin Lỗi Không Phải Mặc Cả}{No Bargaining with Integrity}
\chuhoa{S}{au một lần cô nhận sai}, có vài bạn bắt đầu nhờn. Hễ bị nhắc là lại lôi câu cũ ra: ``Hôm trước cô cũng sai mà.''

Cô phải chặn ngay: ``Đúng. Và cô đã xin lỗi phần cô sai. Nhưng xin lỗi không phải là một loại phiếu giảm giá để từ nay các em được quyền lấn tới. Nếu các em học được từ lời xin lỗi của người lớn chỉ là cách mặc cả để thoát lỗi, thì các em học sai rồi.''

Cả lớp im hẳn, vì đó là khoảnh khắc các em hiểu ra: đạo đức không hoạt động kiểu đổi điểm thưởng.
\end{truyen}

\begin{baihoc}
Xin lỗi là dấu hiệu của nhân cách, không phải sự rút lui khỏi nguyên tắc. Lời xin lỗi có giá trị nhất khi nó vừa sửa lại tổn thương, vừa giữ nguyên tiêu chuẩn cho điều đúng.
\textit{(An apology is a sign of character, not a retreat from principle. The most valuable apology repairs harm while still preserving the standard of what is right.)}
\end{baihoc}

\begin{ghinhoanh}
\textbf{Short English to remember:}

Apology shows strength.

Standards still matter.
\end{ghinhoanh}

\begin{tuhocnhanh}
1. Em có từng nghĩ người lớn xin lỗi là họ yếu đi không? \textit{(Have you ever thought an adult becomes weak by apologizing?)}

2. Em có dám xin lỗi đúng chỗ mà không trốn hậu quả không? \textit{(Can you apologize honestly without escaping the consequence?)}

3. Khi em xin lỗi, em thường mong được hết trách nhiệm luôn hay chỉ mong sửa lại phần mình làm sai? \textit{(When you apologize, do you usually hope to escape all responsibility or to repair the specific part you did wrong?)}
\end{tuhocnhanh}

\begin{tongketchuong}
Một lời xin lỗi tốt không hủy kỷ luật. Nó làm kỷ luật bớt mùi quyền lực và gần hơn với công lý. Trong giáo dục, xin lỗi đúng lúc là một bài dạy sống động hơn rất nhiều bài đạo đức viết trên bảng.
\end{tongketchuong}

\begin{goigiaovien}
Nếu dùng chương này trong bồi dưỡng kỹ năng chủ nhiệm, nên nhấn vào ba điểm: nhận đúng phần sai, sửa đúng chỗ cần sửa, và không biến lời xin lỗi thành màn yếu lòng cảm tính. Học sinh tôn trọng sự mềm đúng chỗ hơn người lớn thường nghĩ.
\end{goigiaovien}

\ngancach